\documentclass[12pt]{extarticle}
\def\activityid{F02-X-v3}\usepackage[letterpaper,margin=.65in,top=.60in,bottom=.65in]{geometry}
\usepackage{fix-cm}
\usepackage[T1]{fontenc}\usepackage[default]{sourcesanspro}
\usepackage{microtype,tikz,array,tabularx,fancyhdr,amsmath,amssymb}
\usepackage[hidelinks]{hyperref}\usetikzlibrary{calc,arrows.meta}
\setlength{\parindent}{0pt}\setlength{\parskip}{7pt}\setlength{\headheight}{0pt}\setlength{\footskip}{26pt}
\setlength{\emergencystretch}{2em}\pagestyle{fancy}\fancyhf{}\renewcommand{\headrulewidth}{0pt}\renewcommand{\footrulewidth}{.4pt}
\fancyfoot[L]{\scriptsize Bellingham Math Circle / Week 2 / \activityid}\fancyfoot[R]{\scriptsize \thepage}
\definecolor{ink}{gray}{.12}\definecolor{soft}{gray}{.95}\color{ink}
\newcommand{\PageTitle}[2]{{\fontsize{10}{12}\selectfont\bfseries WEEK 2\quad TOGGLE LAB\hfill #1\par}\vspace{3pt}{\fontsize{26}{29}\selectfont\bfseries #2\par}\vspace{2pt}}
\newcommand{\NameLine}{{\small Name: \rule{2.65in}{.4pt}\hfill Date: \rule{1.35in}{.4pt}\par}\vspace{4pt}}
\newcommand{\Task}[2]{\par\vspace{3pt}{\large\bfseries #1}\par #2\par}
\newcommand{\Subhead}[1]{\par\vspace{3pt}{\large\bfseries #1}\par}
\newcommand{\RuleBox}[1]{\par\begingroup\setlength{\fboxsep}{8pt}\colorbox{soft}{\parbox{\dimexpr\linewidth-16pt\relax}{#1}}\endgroup\par}
\newcommand{\WriteBox}[2]{\par\textbf{#1}\par\vspace{2pt}\begin{tikzpicture}\draw[black!40,line width=.5pt] (0,0) rectangle (\dimexpr\linewidth-.5pt\relax,#2);\end{tikzpicture}\par}
\newcommand{\StudentFooter}{\vfill{\small Try it with objects. Explain by pointing, talking, or drawing.}\par}
\newcommand{\LampRing}[3]{\begin{tikzpicture}[x=1in,y=1in]
\foreach \i in {1,...,#1}{\coordinate (v\i) at ({90-360*(\i-1)/#1}:#2);}
\foreach \i in {1,...,#1}{\pgfmathtruncatemacro{\j}{mod(\i,#1)+1}\draw[thick] (v\i)--node[midway,fill=white,inner sep=2pt,font=\small]{\i\j}(v\j);}
\foreach \i in {1,...,#1}{\node[circle,draw,fill=white,minimum size=.52in,inner sep=0pt,font=\bfseries] at(v\i){\i};}
\foreach \i in {#3}{\node[circle,draw,fill=ink,text=white,minimum size=.52in,inner sep=0pt,font=\bfseries] at(v\i){\i};}
\end{tikzpicture}}
\newcommand{\Lights}[1]{\begin{tikzpicture}[x=.55in,y=.55in,baseline=-.10in]\foreach \v [count=\i] in {#1}{\ifnum\v=1\fill (\i,0) circle(.16in);\else\draw (\i,0) circle(.16in);\fi}\end{tikzpicture}}
\newcommand{\ClockRing}[2]{\begin{tikzpicture}[x=1in,y=1in]\draw[black!30] (0,0) circle(#2);\pgfmathtruncatemacro{\n}{#1-1}\foreach\i in {0,...,\n}{\node[circle,draw,fill=white,minimum size=.35in,inner sep=0pt] at({90-360*\i/#1}:#2){\i};}\draw[-{Stealth},thick] (-.35,.15) arc(155:25:.38);\node[font=\small] at(0,-.18){clockwise};\end{tikzpicture}}
\newcommand{\Key}[2]{\begin{center}\renewcommand{\arraystretch}{1.55}\begin{tabular}{c|*{#1}{c}}in&#2\end{tabular}\end{center}}


% v3 student pages use one running header and numbered tasks only.
\newcommand{\StudentHeader}[1]{{\fontsize{10}{12}\selectfont\bfseries Week 2 / Toggle lab / #1\par}\vspace{10pt}}
\newcommand{\Problem}[2]{\par\vspace{4pt}\textbf{Problem #1:} #2\par}
\newcommand{\Space}[1]{\begin{tikzpicture}\draw[black!35] (0,0) rectangle (\dimexpr\linewidth-.5pt\relax,#1);\end{tikzpicture}\par}
\newcommand{\Record}[2]{#1\par\Space{#2}}

\begin{document}
\StudentHeader{Extra / Grades 6--7}

\Problem{1}{Place a two-sided counter on every lamp, all OFF. Press an edge to flip both ends. Start over for each target in the table. Record moves that work, or ``not yet.'' A move cannot jump across a gap.}
\begin{center}\begin{tikzpicture}[x=1in,y=1in,every node/.style={circle,draw,fill=white,minimum size=.42in}]\node(A)at(0,.7){A};\node(B)at(-.65,-.4){B};\node(C)at(.65,-.4){C};\draw(A)--(B)--(C)--(A);\node(D)at(2.1,.3){D};\node(E)at(3.1,.3){E};\draw(D)--(E);\end{tikzpicture}\end{center}
\begin{center}\renewcommand{\arraystretch}{2.15}\begin{tabular}{l|p{4in}}ON lamps&Moves, or not yet\\\hline A,C&\\D,E&\\A,D&\\A,C,D,E&\\A,B,C,D&\\A,B,D,E&\\\end{tabular}\end{center}
\Problem{2}{Count the target lamps in each connected piece separately. Sort your targets using those two counts. Record a rule that predicts which targets work.}
\Space{1in}

\newpage
\StudentHeader{Extra / Grades 6--7}

\Problem{3}{Add a new edge CD to the page-1 network. Try the two formerly unsuccessful targets A,D and A,B,C,D again. Record your paths and resulting move lists.}
\Space{1.2in}
\Problem{4}{On the joined network, use paths A--C--D and B--C--D to make target A,B. Try the full list including the repeated CD. Then remove both CD presses and replay it. Compare the final lamps.}
\Space{1.1in}
\Problem{5}{Draw a different network with two connected pieces, at least three lamps in each. Choose one target your rule predicts can be made and one it predicts cannot. Try both and record them on your drawing.}
\Space{2in}

\newpage
\StudentHeader{Extra / Grades 6--7}

\Problem{6}{A tree is a connected network with no closed loop. On this tree, start all OFF and make target A,C. Reset and make A,C,E,F. Find a reduced move list for each: use every edge at most once.}
\begin{center}\begin{tikzpicture}[x=1in,y=1in,every node/.style={circle,draw,fill=white,minimum size=.42in}]\node(A)at(-1.4,.6){A};\node(B)at(0,.6){B};\node(C)at(0,1.35){C};\node(D)at(0,-.3){D};\node(E)at(-1.4,-.3){E};\node(F)at(1.4,-.3){F};\draw(A)--(B)--(C);\draw(B)--(D)--(E);\draw(D)--(F);\end{tikzpicture}\end{center}
\Record{Two move lists:}{.7in}
\Problem{7}{For target A,C,E,F, imagine removing each edge. Count the targets on the named side of the cut. Compare odd/even with whether your move list uses that edge.}
\begin{center}\renewcommand{\arraystretch}{1.8}\begin{tabular}{c|l|c|c}Cut edge&Count on this side&Target count&Edge used?\\\hline AB&A&&\\BC&C&&\\BD&A,B,C&&\\DE&E&&\\DF&F&&\\\end{tabular}\end{center}
\Problem{8}{Choose a different even target. Use cut counts to predict its edge set before testing. Draw or write the prediction, then replay it.}
\Space{.75in}

\end{document}
